% Options for packages loaded elsewhere
\PassOptionsToPackage{unicode}{hyperref}
\PassOptionsToPackage{hyphens}{url}
\PassOptionsToPackage{dvipsnames,svgnames,x11names}{xcolor}
%
\documentclass[
  authoryear,
  preprint]{elsarticle}

\usepackage{amsmath,amssymb}
\usepackage{iftex}
\ifPDFTeX
  \usepackage[T1]{fontenc}
  \usepackage[utf8]{inputenc}
  \usepackage{textcomp} % provide euro and other symbols
\else % if luatex or xetex
  \usepackage{unicode-math}
  \defaultfontfeatures{Scale=MatchLowercase}
  \defaultfontfeatures[\rmfamily]{Ligatures=TeX,Scale=1}
\fi
\usepackage{lmodern}
\ifPDFTeX\else  
    % xetex/luatex font selection
\fi
% Use upquote if available, for straight quotes in verbatim environments
\IfFileExists{upquote.sty}{\usepackage{upquote}}{}
\IfFileExists{microtype.sty}{% use microtype if available
  \usepackage[]{microtype}
  \UseMicrotypeSet[protrusion]{basicmath} % disable protrusion for tt fonts
}{}
\makeatletter
\@ifundefined{KOMAClassName}{% if non-KOMA class
  \IfFileExists{parskip.sty}{%
    \usepackage{parskip}
  }{% else
    \setlength{\parindent}{0pt}
    \setlength{\parskip}{6pt plus 2pt minus 1pt}}
}{% if KOMA class
  \KOMAoptions{parskip=half}}
\makeatother
\usepackage{xcolor}
\setlength{\emergencystretch}{3em} % prevent overfull lines
\setcounter{secnumdepth}{5}
% Make \paragraph and \subparagraph free-standing
\ifx\paragraph\undefined\else
  \let\oldparagraph\paragraph
  \renewcommand{\paragraph}[1]{\oldparagraph{#1}\mbox{}}
\fi
\ifx\subparagraph\undefined\else
  \let\oldsubparagraph\subparagraph
  \renewcommand{\subparagraph}[1]{\oldsubparagraph{#1}\mbox{}}
\fi


\providecommand{\tightlist}{%
  \setlength{\itemsep}{0pt}\setlength{\parskip}{0pt}}\usepackage{longtable,booktabs,array}
\usepackage{calc} % for calculating minipage widths
% Correct order of tables after \paragraph or \subparagraph
\usepackage{etoolbox}
\makeatletter
\patchcmd\longtable{\par}{\if@noskipsec\mbox{}\fi\par}{}{}
\makeatother
% Allow footnotes in longtable head/foot
\IfFileExists{footnotehyper.sty}{\usepackage{footnotehyper}}{\usepackage{footnote}}
\makesavenoteenv{longtable}
\usepackage{graphicx}
\makeatletter
\def\maxwidth{\ifdim\Gin@nat@width>\linewidth\linewidth\else\Gin@nat@width\fi}
\def\maxheight{\ifdim\Gin@nat@height>\textheight\textheight\else\Gin@nat@height\fi}
\makeatother
% Scale images if necessary, so that they will not overflow the page
% margins by default, and it is still possible to overwrite the defaults
% using explicit options in \includegraphics[width, height, ...]{}
\setkeys{Gin}{width=\maxwidth,height=\maxheight,keepaspectratio}
% Set default figure placement to htbp
\makeatletter
\def\fps@figure{htbp}
\makeatother

\makeatletter
\makeatother
\makeatletter
\makeatother
\makeatletter
\@ifpackageloaded{caption}{}{\usepackage{caption}}
\AtBeginDocument{%
\ifdefined\contentsname
  \renewcommand*\contentsname{Table of contents}
\else
  \newcommand\contentsname{Table of contents}
\fi
\ifdefined\listfigurename
  \renewcommand*\listfigurename{List of Figures}
\else
  \newcommand\listfigurename{List of Figures}
\fi
\ifdefined\listtablename
  \renewcommand*\listtablename{List of Tables}
\else
  \newcommand\listtablename{List of Tables}
\fi
\ifdefined\figurename
  \renewcommand*\figurename{Figure}
\else
  \newcommand\figurename{Figure}
\fi
\ifdefined\tablename
  \renewcommand*\tablename{Table}
\else
  \newcommand\tablename{Table}
\fi
}
\@ifpackageloaded{float}{}{\usepackage{float}}
\floatstyle{ruled}
\@ifundefined{c@chapter}{\newfloat{codelisting}{h}{lop}}{\newfloat{codelisting}{h}{lop}[chapter]}
\floatname{codelisting}{Listing}
\newcommand*\listoflistings{\listof{codelisting}{List of Listings}}
\makeatother
\makeatletter
\@ifpackageloaded{caption}{}{\usepackage{caption}}
\@ifpackageloaded{subcaption}{}{\usepackage{subcaption}}
\makeatother
\makeatletter
\@ifpackageloaded{tcolorbox}{}{\usepackage[skins,breakable]{tcolorbox}}
\makeatother
\makeatletter
\@ifundefined{shadecolor}{\definecolor{shadecolor}{rgb}{.97, .97, .97}}
\makeatother
\makeatletter
\makeatother
\makeatletter
\makeatother
\journal{Journal Name}
\ifLuaTeX
  \usepackage{selnolig}  % disable illegal ligatures
\fi
\usepackage[]{natbib}
\bibliographystyle{elsarticle-harv}
\IfFileExists{bookmark.sty}{\usepackage{bookmark}}{\usepackage{hyperref}}
\IfFileExists{xurl.sty}{\usepackage{xurl}}{} % add URL line breaks if available
\urlstyle{same} % disable monospaced font for URLs
\hypersetup{
  pdftitle={Beyond Language Barriers},
  pdfauthor={Laurence-Olivier M. Foisy; Camille Pelletier},
  pdfkeywords={keyword1, keyword2},
  colorlinks=true,
  linkcolor={blue},
  filecolor={Maroon},
  citecolor={Blue},
  urlcolor={Blue},
  pdfcreator={LaTeX via pandoc}}

\setlength{\parindent}{6pt}
\begin{document}

\begin{frontmatter}
\title{Beyond Language Barriers \\\large{Evaluating the Efficacy of
Open-Source LLMs in Analyzing Sentiment in Non-English Textual Data} }
\author[1]{Laurence-Olivier M. Foisy%
\corref{cor1}%
\fnref{fn1}}
 \ead{mail@mfoisy.com} 
\author[1]{Camille Pelletier%
%
}
 \ead{camille.pelletier.12@ulaval.ca} 

\affiliation[1]{organization={Université Laval, Département de science
politique},addressline={2325, rue de
l'Université},city={Québec},postcode={G1V 0A6},postcodesep={}}

\cortext[cor1]{Corresponding author}
\fntext[fn1]{This is the first author footnote.}

        
\begin{abstract}
This is the abstract. Lorem ipsum dolor sit amet, consectetur adipiscing
elit. Vestibulum augue turpis, dictum non malesuada a, volutpat eget
velit. Nam placerat turpis purus, eu tristique ex tincidunt et. Mauris
sed augue eget turpis ultrices tincidunt. Sed et mi in leo porta
egestas. Aliquam non laoreet velit. Nunc quis ex vitae eros aliquet
auctor nec ac libero. Duis laoreet sapien eu mi luctus, in bibendum leo
molestie. Sed hendrerit diam diam, ac dapibus nisl volutpat vitae.
Aliquam bibendum varius libero, eu efficitur justo rutrum at. Sed at
tempus elit.
\end{abstract}





\begin{keyword}
    keyword1 \sep 
    keyword2
\end{keyword}
\end{frontmatter}
    \ifdefined\Shaded\renewenvironment{Shaded}{\begin{tcolorbox}[breakable, interior hidden, borderline west={3pt}{0pt}{shadecolor}, frame hidden, boxrule=0pt, sharp corners, enhanced]}{\end{tcolorbox}}\fi

\hypertarget{introduction}{%
\section{Introduction}\label{introduction}}

The rise of large language models (LLMs) has transformed natural
language processing across domains and applications. While these models
demonstrate impressive capabilities in English, their cross-linguistic
effectiveness seems to remain a critical area of investigation. This
paper addresses a fundamental question in this domain: Can open source
general purpose foundational LLMs accurately evaluate sentiment in
non-English texts?

Recent research has uncovered a compelling phenomenon: prompting in
English often improves performance in cross-linguistic tasks, even when
the target language is non-English. Jiang et al.~(2019) documented
increased accuracy in knowledge extraction when using English prompts
compared to diversified manual prompts in other languages. This finding
was corroborated by Lin et al.~(2021), who observed that English prompts
outperform multilingual configurations across natural language
understanding and machine translation tasks in 30 languages, even
surpassing established baselines like GPT-3.

This pattern extends beyond isolated studies. Muennighoff et al.~(2023),
Nie et al.~(2024), have all reported significant performance
gains---measured by accuracy, F1 scores, and BLEU metrics---when models
are guided with English instructions, regardless of the target language.
These findings suggest a fundamental characteristic of current LLM
architectures that merits deeper investigation.

The primary purpose of this research is to determine whether untrained
foundational models are adequate for analyzing textual sentiment in
non-English contexts. While specialized models like BERT and its
variants have been fine-tuned specifically for sentiment analysis tasks
and perform effectively, they often require technical expertise to
implement and adapt. By contrast, using non-fine-tuned foundational
models could provide a more accessible solution for non-technical social
science researchers who may lack the computational resources or
programming knowledge necessary for model customization. This
accessibility consideration is particularly relevant as AI tools become
increasingly integrated into research methodologies across disciplines.

Sentiment analysis offers an interesting test case for examining this
phenomenon, as it requires nuanced language understanding beyond simple
translation. While proprietary models have received substantial
attention in multilingual contexts, open source LLMs present both unique
challenges and opportunities. They offer greater transparency,
adaptability, and accessibility for global researchers and
practitioners, yet their cross-linguistic capabilities remain less
comprehensively documented. This study seeks to evaluate the performance
of leading open source foundational LLMs in sentiment analysis with
french texts. We examine how prompt language affects performance,
identify cross-linguistic patterns, and analyze model behavior across
french and english. Our findings seeks to contribute to a deeper
understanding of LLMs capabilities and limitations in multilingual
contexts and offer practical guidelines for optimizing their deployment
in diverse language settings.

\hypertarget{literature-review}{%
\section{Literature Review}\label{literature-review}}

\hypertarget{specialized-models-for-multilingual-sentiment-analysis}{%
\subsection{Specialized Models for Multilingual Sentiment
Analysis}\label{specialized-models-for-multilingual-sentiment-analysis}}

Previous research consistently demonstrates the effectiveness of
specialized models for sentiment analysis in non-English languages.
These models, based on various transformer architectures, have achieved
impressive performance across diverse linguistic contexts.
Transformer-based architectures have shown particularly strong results,
with accuracy rates often exceeding 85\% across various Indo-European
and Semitic languages (Bhowmick \& Jana, 2021; Michailidis, 2024).
Researchers have reported up to 95\% accuracy for Bengali sentiment
analysis (Bhowmick \& Jana, 2021), while others achieved 96\% accuracy
for Greek sentiment analysis (Michailidis, 2024). For Arabic, Hameed et
al.~(2023) documented 96.442\% accuracy using specialized models,
demonstrating the effectiveness of language-specific fine-tuning.

These language-specific adaptations consistently outperform general
multilingual approaches. ElJundi et al.~(2019) and Ahmadian et
al.~(2024) developed specialized models for Arabic and Indonesian
respectively, both reporting significant improvements over generic
multilingual models. This pattern holds across diverse language
families, with specialized models for South Asian, East Asian, and
Semitic languages demonstrating strong performance within their target
domains (Gaikwad et al., 2023; Saeed et al., 2024). Even for
lower-resource languages, fine-tuned models have shown promise. Tela et
al.~(2020) demonstrated competitive performance for Tigrinya, while
Eusha et al.~(2024) reported encouraging results for Tamil and Tulu
using specialized transformer-based architectures. Salahudeen et
al.~(2023) documented F1-scores between 62\% and 83\% for various
African languages, noting that performance improved with increased
training data availability.

The literature also shows progress in handling code-mixed and
transliterated text, particularly relevant for many South Asian language
contexts. Nazir et al.~(2025) reported significant improvements in
analyzing code-mixed Roman Urdu and Roman Punjabi, while Ipa et
al.~(2024) introduced specialized models for Bangla-English code-mixed
content. For multilingual applications, cross-lingual models demonstrate
strong transfer capabilities. Kumar and Albuquerque (2021) successfully
employed advanced techniques for zero-shot transfer from English to
Hindi, while Pan et al.~(2023) highlighted the effectiveness of
cross-lingual transfer learning for financial sentiment analysis in
Spanish.

\hypertarget{the-resource-performance-tradeoff}{%
\subsection{The Resource-Performance
Tradeoff}\label{the-resource-performance-tradeoff}}

Despite their impressive results, specialized sentiment analysis models
present implementation challenges. As Přibáň et al.~(2024) discuss,
there exists a clear tradeoff between performance and efficiency, with
advanced models requiring computational resources. Language-specific
models typically contain millions of parameters (Hameed et al., 2023),
while larger multilingual models can expand to hundreds of millions or
even billions of parameters (Barbieri et al., 2021; Vileikytė et al.,
2024).

The implementation of these specialized models requires not only
significant computational infrastructure but also technical expertise in
model selection, fine-tuning, and deployment. Kotelnikova et al.~(2023)
demonstrated meaningful improvements after fine-tuning on
domain-specific Russian data, but such fine-tuning processes demand
substantial ML engineering knowledge. For many social science
researchers and practitioners without technical backgrounds, these
requirements present barriers to adoption.

Additionally, the literature emphasizes the importance of high-quality,
domain-specific training data. Salahudeen et al.~(2023) noted better
performance for Nigerian languages with larger available datasets, while
Krasitskii et al.~(2024) highlighted the challenges of limited data
availability for linguistically diverse languages like Finnish,
Hungarian, and Bulgarian. These data requirements further complicate
implementation for researchers working with low-resource languages or
specialized domains.

\hypertarget{the-potential-of-general-purpose-foundational-llms}{%
\subsection{The Potential of General-Purpose Foundational
LLMs}\label{the-potential-of-general-purpose-foundational-llms}}

While specialized models dominate the literature on multilingual
sentiment analysis, relatively little attention has been given to the
capabilities of general-purpose foundational LLMs without task-specific
fine-tuning. This represents a gap in the research, particularly given
the findings of Jiang et al.~(2019), Lin et al.~(2021), and others
regarding the unexpected effectiveness of English-prompted models on
cross-linguistic tasks.

Recent studies by Venerito and Iannone (2024) and Buscemi and Proverbio
(2024) have begun to explore the potential of general-purpose LLMs like
Mistral-7B and LLaMA2 for sentiment analysis in Italian and across
multiple languages, respectively. However, these studies remain
preliminary, and comprehensive evaluations of open-source foundational
LLMs across diverse language remain limited. This gap is particularly
notable given the practical advantages general-purpose LLMs might offer.
As highlighted by Plagwitz et al.~(2024) in their work on German
clinical texts, leveraging foundational models without extensive
fine-tuning could reduce implementation barriers. Similarly, Rusnachenko
et al.~(2024) observed promising results using general instruction-tuned
models for sentiment analysis in Russian and English, suggesting the
potential of generalist models for cross-linguistic applications.

The limited existing research on general-purpose LLMs for multilingual
sentiment analysis reveals several patterns worth investigating.
Vileikytė et al.~(2024) found that language-specific prompting
strategies significantly impact performance for Lithuanian sentiment
analysis, while Hämäläinen et al.~(2022) documented varying
effectiveness across Germanic and Romance languages depending on prompt
design and model architecture.

\hypertarget{research-question}{%
\section{Research Question}\label{research-question}}

Despite the extensive literature on specialized sentiment analysis
models, a critical question remains largely unexplored: Can open-source
general-purpose foundational LLMs accurately evaluate sentiment in
non-English texts without task-specific fine-tuning? This question has
significant implications for research accessibility, as it could
potentially democratize sophisticated NLP capabilities for non-technical
researchers across disciplines.

The phenomenon documented by Muennighoff et al.~(2023), Nie et
al.~(2024), and others---that English prompting often improves
cross-linguistic performance---suggests that foundational LLMs may
possess unexpectedly robust multilingual capabilities even without
specialized training. This aligns with observations from Rusnachenko et
al.~(2024) and Venerito and Iannone (2024), who found promising results
with general-purpose models in Russian and Italian contexts,
respectively. Our study addresses this research gap by evaluating
leading open-source foundational LLMs with french texts.

\hypertarget{data-methods}{%
\section{Data \& Methods}\label{data-methods}}

\hypertarget{results}{%
\section{Results}\label{results}}

\hypertarget{discussion}{%
\section{Discussion}\label{discussion}}


\renewcommand\refname{References}
  \bibliography{bibliography.bib}


\end{document}
